\section{Fairness: a structured research draft}
A proposed structural interpretation is that no player position is privileged
when the automorphism group acts transitively on $I$:
\[
 \forall i,j\in I\quad\exists g\in\Aut(G)\quad q(g)(i)=j.
\]
This is already Ham's label-independent symmetry condition
\citep[Definition~5.4]{ham2021}, rather than a new theorem about fairness.
Adopting it as a normative definition requires an explicit account of what
players regard as equivalent opportunities and of how utilities are compared.

\begin{question}
Which interpretations of indifference to positions preserve the order of
opponents? Which permit strategy relabeling? Do they recover existing
symmetry notions, or define distinct ones?
\end{question}
\begin{question}
Under which specified distributions does structural symmetry imply equality
of expected payoffs? Separate uniform profiles, common independent mixed
strategies, invariant correlated distributions, and equilibrium selection.
\end{question}
Equality under every distribution on $S$ is a particularly strong requirement:
applying it to point masses yields equality at every pure profile. It therefore
cannot serve as a general synonym for structural symmetry.
The converse follows by linearity of expectation.


Von Neumann and Morgenstern's Subsection~28.2.2 defines fairness by requiring
the coalition characteristic function to depend only on coalition cardinality
\citep{vnm1953}. This belongs to their zero-sum framework and does not itself
select an equilibrium or allocation. Relating it to general-sum normal-form
opportunity equivalence requires a specified coalition-value construction
and a separate argument. We therefore retain that relationship as a research
question.
